\documentclass[10pt,a4paper]{amsart}
\usepackage{amsmath,amssymb,amsthm,graphicx}
\usepackage{verbatim}
\newenvironment{reference}{\comment}{\endcomment}
%\newenvironment{reference}{}{}
\newenvironment{slogan}{\comment}{\endcomment}
\newenvironment{history}{\comment}{\endcomment}

% For commutative diagrams we use Xy-pic
\usepackage[all]{xy}

% We use 2cell for 2-commutative diagrams.
\xyoption{2cell}
% Theorem environments.
%
\theoremstyle{plain}
\newtheorem{theorem}[subsection]{Theorem}
\newtheorem{proposition}[subsection]{Proposition}
\newtheorem{lemma}[subsection]{Lemma}

\theoremstyle{definition}
\newtheorem{definition}[subsection]{Definition}
\newtheorem{example}[subsection]{Example}
\newtheorem{exercise}[subsection]{Exercise}
\newtheorem{situation}[subsection]{Situation}

\theoremstyle{remark}
\newtheorem{remark}[subsection]{Remark}
\newtheorem{remarks}[subsection]{Remarks}

\numberwithin{equation}{subsection}

% Macros
%
\def\lim{\mathop{\mathrm{lim}}\nolimits}
\def\colim{\mathop{\mathrm{colim}}\nolimits}
\def\Spec{\mathop{\mathrm{Spec}}}
\def\Hom{\mathop{\mathrm{Hom}}\nolimits}
\def\Ext{\mathop{\mathrm{Ext}}\nolimits}
\def\SheafHom{\mathop{\mathcal{H}\!\mathit{om}}\nolimits}
\def\SheafExt{\mathop{\mathcal{E}\!\mathit{xt}}\nolimits}
\def\Sch{\mathit{Sch}}
\def\Mor{\mathop{\mathrm{Mor}}\nolimits}
\def\Ob{\mathop{\mathrm{Ob}}\nolimits}
\def\Sh{\mathop{\mathit{Sh}}\nolimits}
\def\NL{\mathop{N\!L}\nolimits}
\def\CH{\mathop{\mathrm{CH}}\nolimits}
\def\proetale{{pro\text{-}\acute{e}tale}}
\def\etale{{\acute{e}tale}}
\def\QCoh{\mathit{QCoh}}
\def\Ker{\mathop{\mathrm{Ker}}}
\def\Im{\mathop{\mathrm{Im}}}
\def\Coker{\mathop{\mathrm{Coker}}}
\def\Coim{\mathop{\mathrm{Coim}}}

% Boxtimes
%
\DeclareMathSymbol{\boxtimes}{\mathbin}{AMSa}{"02}

%
% Macros for moduli stacks/spaces
%
\def\QCohstack{\mathcal{QC}\!\mathit{oh}}
\def\Cohstack{\mathcal{C}\!\mathit{oh}}
\def\Spacesstack{\mathcal{S}\!\mathit{paces}}
\def\Quotfunctor{\mathrm{Quot}}
\def\Hilbfunctor{\mathrm{Hilb}}
\def\Curvesstack{\mathcal{C}\!\mathit{urves}}
\def\Polarizedstack{\mathcal{P}\!\mathit{olarized}}
\def\Complexesstack{\mathcal{C}\!\mathit{omplexes}}
% \Pic is the operator that assigns to X its picard group, usage \Pic(X)
% \Picardstack_{X/B} denotes the Picard stack of X over B
% \Picardfunctor_{X/B} denotes the Picard functor of X over B
\def\Pic{\mathop{\mathrm{Pic}}\nolimits}
\def\Picardstack{\mathcal{P}\!\mathit{ic}}
\def\Picardfunctor{\mathrm{Pic}}
\def\Deformationcategory{\mathcal{D}\!\mathit{ef}}
\usepackage{cleveref}
\crefname{originallemma}{original lemma}{original lemmas}
\Crefname{originallemma}{Original lemma}{Original lemmas}
\makeatletter
\newlabel{algebra-lemma-Mittag-Leffler}{{03CA}{source}{}{sourcepointer-1}{}}
\newlabel{algebra-lemma-Mittag-Leffler@cref}{{[originallemma][0][]03CA}{[1][0][]source}{}{}{}}
\makeatother
\crefname{originallemma}{original lemma}{original lemmas}
\Crefname{originallemma}{Original lemma}{Original lemmas}
\makeatletter
\newlabel{algebra-lemma-characterize-local-ring}{{00E9}{source}{}{sourcepointer-2}{}}
\newlabel{algebra-lemma-characterize-local-ring@cref}{{[originallemma][0][]00E9}{[1][0][]source}{}{}{}}
\makeatother
\crefname{originallemma}{original lemma}{original lemmas}
\Crefname{originallemma}{Original lemma}{Original lemmas}
\makeatletter
\newlabel{algebra-lemma-hathat}{{0318}{source}{}{sourcepointer-3}{}}
\newlabel{algebra-lemma-hathat@cref}{{[originallemma][0][]0318}{[1][0][]source}{}{}{}}
\makeatother
\crefname{originalsection}{original section}{original sections}
\Crefname{originalsection}{Original section}{Original sections}
\makeatletter
\newlabel{more-algebra-section-topological-ring}{{07E7}{source}{}{sourcepointer-4}{}}
\newlabel{more-algebra-section-topological-ring@cref}{{[originalsection][0][]07E7}{[1][0][]source}{}{}{}}
\makeatother
\title{A local ring whose m-adic completion is local but not complete and whose maximal ideal is strictly larger than the expanded one}
\author{The Stacks Project Authors}
\date{stacks-project:examples.tex:local-snapshot-20261012; GNU FDL 1.2 or later}
\begin{document}
\maketitle
\noindent\textit{Editorial note.} All delivered mathematical text is the author's own, copied from the pinned Stacks Project chapter file examples.tex; no statement or proof step was rewritten, reordered or removed.\par\medskip
\noindent\textit{Editorial note.} The counted proof is the author's complete argument: The counted proof adds only the localization bridge ('the completion of the localization R\_m is equal to the completion of R'), the author's own one-sentence assertion; the three verifications live in the context span and are complete there.\par\medskip
\noindent\textit{Editorial note.} The label lemma-noncomplete-completion names the counted statement, and the proof environment immediately following it in the source is delivered verbatim as its proof.\par\medskip
\noindent\textit{Editorial note.} The counted statement is carried by the author's own lemma environment in the chapter source.\par\medskip
\noindent\textit{Allowed standard prerequisites (editorial).}\begin{itemize}
\item Context span 525-651: for R = k[x\_1, x\_2, ...] with m = (x\_1, x\_2, ...), the completion R\textasciicircum{} = lim R/m\textasciicircum{}n is local with maximal ideal m' = Ker(R\textasciicircum{} -> R/m); the element f = x\_1 + x\_2\textasciicircum{}2 + x\_3\textasciicircum{}3 + ... lies in m' but not in m R\textasciicircum{} (falsifying (1)); Algebra lemma-hathat with n=1 then falsifies (3); the Mittag-Leffler exact sequences reduce m'-completeness to K\_n = (m')\textasciicircum{}n, and the explicit z = sum x\_i\textasciicircum{}d with d prime to char(k) (vanishing locus of partial derivatives one point vs Leibniz rule) shows K\_2 != (m')\textasciicircum{}2, falsifying (2) - the whole argument from an unpublished note of de Smit and Lenstra, complete in the section
\item Algebra, lemma-characterize-local-ring, lemma-hathat, lemma-Mittag-Leffler
\end{itemize}\medskip
\ifdefined\hypertarget\hypertarget{sourcepointer-1}{}\fi
\noindent\textit{Source reference.} Original lemma 03CA: sources/stacks-chow/tags.txt, source lines 4377--4377. Prerequisite pointer invoked at 734; it names a standard fact of a sibling chapter that the counted argument uses, and every step of the counted proof built on it is retained verbatim in the delivered body.\par\medskip
\ifdefined\hypertarget\hypertarget{sourcepointer-2}{}\fi
\noindent\textit{Source reference.} Original lemma 00E9: sources/stacks-chow/tags.txt, source lines 526--526. Prerequisite pointer invoked at 682; it names a standard fact of a sibling chapter that the counted argument uses, and every step of the counted proof built on it is retained verbatim in the delivered body.\par\medskip
\ifdefined\hypertarget\hypertarget{sourcepointer-3}{}\fi
\noindent\textit{Source reference.} Original lemma 0318: sources/stacks-chow/tags.txt, source lines 3975--3975. Prerequisite pointer invoked at 713; it names a standard fact of a sibling chapter that the counted argument uses, and every step of the counted proof built on it is retained verbatim in the delivered body.\par\medskip
\ifdefined\hypertarget\hypertarget{sourcepointer-4}{}\fi
\noindent\textit{Source reference.} Original section 07E7: sources/stacks-chow/tags.txt, source lines 9437--9437. Prerequisite pointer invoked at 685; it names a standard fact of a sibling chapter that the counted argument uses, and every step of the counted proof built on it is retained verbatim in the delivered body.\par\medskip

\noindent
Let $R$ be a ring and let $\mathfrak m$ be a maximal ideal. Consider the
completion
$$
R^\wedge = \lim R/\mathfrak m^n.
$$
Note that $R^\wedge$ is a local ring with maximal ideal
$\mathfrak m' = \Ker(R^\wedge \to R/\mathfrak m)$.
Namely, if $x = (x_n) \in R^\wedge$ is not in $\mathfrak m'$, then
$y = (x_n^{-1}) \in R^\wedge$ satisfies $xy = 1$, whence $R^\wedge$ is local by
Algebra, Lemma \ref{algebra-lemma-characterize-local-ring}. Now it is
always true that $R^\wedge$ complete in its limit topology (see the
discussion in
More on Algebra, Section \ref{more-algebra-section-topological-ring}).
But beyond that, we have the following questions:
\begin{enumerate}
\item Is it true that $\mathfrak m R^\wedge = \mathfrak m'$?
\item Is $R^\wedge$ viewed as an $R^\wedge$-module $\mathfrak m'$-adically
complete?
\item Is $R^\wedge$ viewed as an $R$-module $\mathfrak m$-adically complete?
\end{enumerate}
It turns out that these questions all have a negative answer.

\medskip\noindent
Let $k$ be a field, $R = k[x_1, x_2, x_3, \ldots]$, and
$\mathfrak m = (x_1, x_2, x_3, \ldots)$.
We will think of an element $f$ of $R^\wedge$ as a (possibly) infinite sum
$$
f = \sum a_I x^I
$$
(using multi-index notation) such that for each $d \geq 0$ there
are only finitely many nonzero $a_I$ for $|I| = d$. The maximal
ideal $\mathfrak m' \subset R^\wedge$ is the collection of $f$ with
zero constant term. In particular, the element
$$
f = x_1 + x_2^2 + x_3^3 + \ldots
$$
is in $\mathfrak m'$ but not in $\mathfrak m R^\wedge$ which
shows that (1) is false in this example. However, if (1) is
false, then (3) is necessarily false because
$\mathfrak m' = \Ker(R^\wedge \to R/\mathfrak m)$
and we can apply
Algebra, Lemma \ref{algebra-lemma-hathat} with $n = 1$.

\medskip\noindent
To finish we prove that $R^\wedge$ is not $\mathfrak m'$-adically complete.
For $n \geq 1$ let $K_n = \Ker(R^\wedge \to R/\mathfrak m^n)$. Then
we have short exact sequences
$$
0 \to K_n/(\mathfrak m')^n \to R^\wedge/(\mathfrak m')^n \to
R/\mathfrak m^n \to 0
$$
The projection map $R^\wedge \to R/\mathfrak m^{n + 1}$ sends
$(\mathfrak m')^n$ onto $\mathfrak m^n/\mathfrak m^{n + 1}$.
It follows that $K_{n + 1} \to K_n/(\mathfrak m')^n$ is surjective.
Hence the inverse system $\left(K_n/(\mathfrak m')^n\right)$
has surjective transition maps and
taking inverse limits we obtain an exact sequence
$$
0 \to \lim K_n/(\mathfrak m')^n \to
\lim R^\wedge/(\mathfrak m')^n \to
\lim R/\mathfrak m^n \to 0
$$
by Algebra, Lemma \ref{algebra-lemma-Mittag-Leffler}.
Thus we see that $R^\wedge$ is complete with respect to $\mathfrak m'$
if and only if $K_n = (\mathfrak m')^n$ for all $n \geq 1$.

\medskip\noindent
To show that $R^\wedge$ is not $\mathfrak m'$-adically complete
in our example we show that $K_2 = \Ker(R^\wedge \to R/\mathfrak m^2)$
is not equal to $(\mathfrak m')^2$.
Note that an element of $(\mathfrak m')^2$
can be written as a finite sum
\begin{equation}
\label{equation-sum}
\sum\nolimits_{i = 1, \ldots, t} f_i g_i
\end{equation}
with $f_i, g_i \in R^\wedge$ having vanishing constant terms.
To get an example we are going to choose an $z \in K_2$
of the form
$$
z = z_1 + z_2 + z_3 + \ldots
$$
with the following properties
\begin{enumerate}
\item there exist sequences $1 < d_1 < d_2 < d_3 < \ldots $ and
$0 < n_1 < n_2 < n_3 < \ldots$ such that
$z_i \in k[x_{n_i}, x_{n_i + 1}, \ldots, x_{n_{i + 1} - 1}]$
homogeneous of degree $d_i$, and
\item in the ring $k[[x_{n_i}, x_{n_i + 1}, \ldots, x_{n_{i + 1} - 1}]]$
the element $z_i$ cannot be written as a sum (\ref{equation-sum})
with $t \leq i$.
\end{enumerate}
Clearly this implies that $z$ is not in $(\mathfrak m')^2$
because the image of the relation (\ref{equation-sum}) in the
ring $k[[x_{n_i}, x_{n_i + 1}, \ldots, x_{n_{i + 1} - 1}]]$
for $i$ large enough would produce a contradiction. Hence it suffices
to prove that for all $t > 0$ there exists a $d \gg 0$ and an integer
$n$ such that we can find an homogeneous element
$z \in k[x_1, \ldots, x_n]$ of degree $d$ which cannot be written as
a sum (\ref{equation-sum}) for the given $t$ in $k[[x_1, \ldots, x_n]]$.
Take $n > 2t$ and any $d > 1$ prime to the characteristic of $k$ and
set $z = \sum_{i = 1, \ldots, n} x_i^d$. Then the vanishing locus
of the ideal
$$
(\frac{\partial z}{\partial x_1}, \ldots, \frac{\partial z}{\partial x_n})
=
(dx_1^{d - 1}, \ldots, dx_n^{d - 1})
$$
consists of one point. On the other hand,
$$
\frac{\partial ( \sum\nolimits_{i = 1, \ldots, t} f_i g_i ) }{\partial x_j}
\in (f_1, \ldots, f_t, g_1, \ldots, g_t)
$$
by the Leibniz rule and hence the vanishing locus of these derivatives
contains at least
$$
V(f_1, \ldots, f_t, g_1, \ldots, g_t) \subset
\Spec(k[[x_1, \ldots, x_n]]).
$$
Hence this is a contradiction as the dimension of
$V(f_1, \ldots, f_t, g_1, \ldots, g_t)$ is at least $n - 2t \geq 1$.

\begin{lemma}
\label{lemma-noncomplete-completion}
There exists a local ring $R$ and a maximal ideal $\mathfrak m$ such that
the completion $R^\wedge$ of $R$ with respect to $\mathfrak m$ has the
following properties
\begin{enumerate}
\item $R^\wedge$ is local, but its maximal ideal is not equal to
$\mathfrak m R^\wedge$,
\item $R^\wedge$ is not a complete local ring, and
\item $R^\wedge$ is not $\mathfrak m$-adically complete as an $R$-module.
\end{enumerate}
\end{lemma}

\begin{proof}
This follows from the discussion above as (with $R = k[x_1, x_2, x_3, \ldots]$)
the completion of the localization $R_{\mathfrak m}$ is equal to the
completion of $R$.
\end{proof}

\end{document}
